\section{مقدمة}

تُعد النمذجة الديناميكية من الركائز الأساسية في دراسة الأنظمة الروبوتية، إذ تهدف إلى توصيف العلاقة بين العزوم المطبقة على مفاصل الذراع الروبوتية والحركة الناتجة عنها. وعلى عكس النمذجة الحركية التي تركز على تحديد موضع واتجاه النهاية الفعالة اعتمادًا على المتغيرات المفصلية، فإن النمذجة الديناميكية تأخذ بعين الاعتبار الخصائص الفيزيائية للنظام، مثل الكتل، وعزوم القصور الذاتي، وتأثيرات الجاذبية، بالإضافة إلى التفاعلات المتبادلة بين الوصلات المختلفة.

بعد الانتهاء من توصيف البنية الحركية للذراع الروبوتية في الفصل السابق، والتي شملت اشتقاق الحركيات الأمامية والعكسية باستخدام معاملات دينافيت--هارتنبرغ، يصبح من الضروري الانتقال إلى دراسة السلوك الديناميكي للنظام. إذ تتيح هذه الدراسة تحديد العزوم اللازمة عند المفاصل من أجل تحقيق حركة معينة أو المحافظة على وضعية محددة للذراع، وهو ما يُعد خطوة أساسية قبل الانتقال إلى تصميم أنظمة التحكم وتحليل استجابة النظام.

في هذا الفصل، سيتم اشتقاق النموذج الديناميكي للذراع الروبوتية ذات خمس درجات حرية باستخدام صياغة لاغرانج، وذلك من خلال حساب طاقة الحركة وطاقة الوضع للنظام، ثم تطبيق معادلات لاغرانج للحصول على معادلات الحركة. كما سيتم استخراج مصفوفة العطالة، ومركبات الجاذبية، وقصور المشغلات، ثم استخدامها للوصول إلى معادلات العزوم النهائية للمفاصل.

ونظرًا لطول المعادلات الناتجة بعد التعويض الكامل في النموذج الديناميكي، سيتم عرض النموذج بصيغته المصفوفية العامة داخل متن الفصل، مع إرفاق المعادلات التفصيلية في ملف خارجي مستقل. كما سيتم اعتماد بعض التبسيطات المناسبة لطبيعة النظام، مثل إهمال قوى كوريوليس والقوى النابذة المركزية في هذه المرحلة، وذلك بهدف التركيز على تأثير العطالة والجاذبية في حساب عزوم المفاصل.

بالإضافة إلى الاشتقاق الرياضي، سيتم استخدام أدوات المحاكاة مثل \LR{SolidWorks} و\LR{Simscape} \LR{Multibody} لبناء نموذج محاكاة فيزيائي للذراع الروبوتية داخل بيئة \LR{MATLAB}/\LR{Simulink}. ويساعد هذا النموذج على تمثيل البنية الميكانيكية للذراع بصورة قريبة من التصميم الحقيقي، ومتابعة حركة الوصلات والمفاصل ضمن بيئة المحاكاة.

يُشكّل هذا الفصل الأساس النظري والعددي لفهم السلوك الديناميكي للذراع الروبوتية، كما يمهّد للمراحل اللاحقة من العمل، سواء في حساب العزوم المطلوبة، أو اختبار حركة الذراع، أو ربط النموذج مع خوارزميات التحكم.